O coeficiente de volume para um sólido isotrópico é $\beta\approx3\alpha$, enquanto o coeficiente de área é $\gamma\approx2\alpha$. Usar $\beta$ para estimar mudança de área produziria um erro de ordem relativa
\[
\frac{\beta-\gamma}{\gamma}=\frac{3\alpha-2\alpha}{2\alpha}=\frac12 = 50\%.
\]
Esse erro não tende a zero com $\Delta T$ porque o coeficiente usado está incorreto para área. Portanto não existe uma diferença de temperatura finita para a qual esse erro seja apenas 1%; o uso do coeficiente de volume para calcular mudança de área é incompatível em termos de coeficiente básico. 